% ============================================================
\part{Mit Textdateien und Skripten arbeiten}
% ============================================================

\chapter{Wir wollen eine Konfigurationsdatei bearbeiten}

\kapitelziel{
Du erstellst die Sensor-Konfiguration für robotik-uno-q
in YAML und kannst sie im Terminal bearbeiten.
}
\kapitelstruktur

\section{Situation}

\texttt{sensor\_bridge.py} hat zurzeit den Schwellenwert
für den Sicherheitsabstand hart im Code stehen:
\texttt{MIN\_ABSTAND = 0.20}. Wenn du ihn ändern willst,
musst du jedes Mal den Python-Code öffnen.

Das ist unprofessionell -- und gefährlich.
Ein Konfigurationsfehler soll nie den Code berühren müssen.

\begin{aufgabeBox}
Erstelle \texttt{stufe2\_ros2/config/sensor\_config.yaml}
mit den HC-SR04-Parametern des Roboters und
binde die Konfiguration in \texttt{sensor\_bridge.py} ein.
\end{aufgabeBox}

\begin{problemBox}
YAML ist einrückungsbasiert -- ein einzelnes falsches Leerzeichen
macht die gesamte Datei ungültig. Und Tab-Einrückung
ist in YAML verboten.
\end{problemBox}

\begin{haubeBox}
\textbf{Warum Konfiguration von Code trennen?}

Derselbe Code soll auf verschiedenen Robotern laufen.
Roboter~A hat Sensoren mit 2\,m Reichweite, Roboter~B mit 4\,m.
Wenn die Reichweite im Code steht, brauchst du für jeden
Roboter eine eigene Code-Version.
Wenn sie in einer Konfigurationsdatei steht, reicht eine
andere \texttt{.yaml}-Datei -- der Code bleibt identisch.

Das ist das Prinzip hinter ROS2-Parametern,
Docker-Compose-Umgebungsvariablen und
Linux-Systemkonfigurationsdateien in \texttt{/etc/}.
Überall dasselbe Muster: Code = Logik, Konfiguration = Daten.
\end{haubeBox}

\section{Die Sensor-Konfiguration}

\begin{lstlisting}[style=yaml]
# stufe2_ros2/config/sensor_config.yaml
# Konfiguration fuer HC-SR04 Ultraschall-Sensoren
# Hardware: robotik-uno-q mit Arduino UNO Q / UNO R3

sensoren:
  ultraschall_rechts:
    trigger_pin: 3     # D3 am Arduino (aus hinderniserkennung.ino)
    echo_pin: 4        # D4 am Arduino
    max_abstand_m: 4.0
    einheit: "meter"

  ultraschall_links:
    trigger_pin: 9     # D9 am Arduino
    echo_pin: 10       # D10 am Arduino
    max_abstand_m: 4.0
    einheit: "meter"

sicherheit:
  min_abstand_m: 0.25    # Stopp-Schwellenwert
  warn_abstand_m: 0.50   # Warnungs-Schwellenwert

serial:
  baudrate: 9600
  timeout_s: 1.0
  encoding: "utf-8"
\end{lstlisting}

\begin{loesungBox}
\begin{lstlisting}[style=bash]
cd ~/robotik-uno-q/stufe2_ros2
nano config/sensor_config.yaml
# Inhalt eingeben (siehe oben)
# Strg+O Enter (speichern), Strg+X (beenden)

# Validieren:
python3 -c "import yaml; print(yaml.safe_load(open('config/sensor_config.yaml')))"
\end{lstlisting}

Einbinden in \texttt{sensor\_bridge.py}:
\begin{lstlisting}[style=python]
#!/usr/bin/env python3
# stufe2_ros2/src/sensor_bridge.py -- Version 0.2
import yaml

with open('../config/sensor_config.yaml') as f:
    config = yaml.safe_load(f)

min_abstand = config['sicherheit']['min_abstand_m']
baudrate    = config['serial']['baudrate']

print(f"[sensor_bridge] Konfiguration geladen.")
print(f"  Min-Abstand: {min_abstand} m")
print(f"  Baudrate:    {baudrate}")
\end{lstlisting}
\end{loesungBox}

\begin{projektBox}
\textbf{Bisher:} \texttt{src/sensor\_bridge.py} mit hartem Schwellenwert.\\
\textbf{Heute:} \texttt{config/sensor\_config.yaml} erstellt,
Konfiguration ausgelagert, YAML-Syntax gelernt.\\
\textbf{Nächstes Kapitel:} Startskript schreiben.
\end{projektBox}

\begin{merksatzBox}
YAML-Regeln: Einrückung nur mit Leerzeichen (nie Tabs).
Doppelpunkt mit Leerzeichen danach: \texttt{key: value}.
Listen mit Bindestrich: \texttt{- item}.
Kommentare mit \texttt{\#}.
\end{merksatzBox}

\begin{fehlerBox}
\begin{itemize}
  \item \textbf{Tab-Einrückung:} \texttt{yaml.scanner.ScannerError}
        -- oft durch Copy-Paste aus einem Editor mit Tabs.
  \item \textbf{Fehlender Doppelpunkt-Abstand:}
        \texttt{key:value} ist falsch, \texttt{key: value} ist richtig.
\end{itemize}
\end{fehlerBox}

\begin{robotikBox}
ROS2-Parameter werden in YAML-Dateien gespeichert und beim
Node-Start übergeben:
\begin{lstlisting}[style=bash]
ros2 run robotik_uno_q sensor_node \
    --ros-args --params-file config/sensor_config.yaml
\end{lstlisting}
Die Datei, die du gerade erstellt hast, ist in
Band~2 bereits die richtige Grundlage.
\end{robotikBox}

% ── Kapitel 9 ────────────────────────────────────────────────
\chapter{Wir wollen mehrere Befehle automatisch ausführen}

\kapitelziel{
Du schreibst ein Shell-Skript, das den robotik-uno-q-Stack
vollständig startet -- ein Befehl für alles.
}
\kapitelstruktur

\section{Situation}

Um \texttt{sensor\_bridge.py} zu starten, musst du jedes Mal:
\begin{enumerate}
  \item In das richtige Verzeichnis wechseln
  \item ROS2 laden (falls installiert)
  \item Pyserial prüfen
  \item Das Skript starten
\end{enumerate}

Vier Befehle. Jeden Tag. Immer in derselben Reihenfolge.
Das ist eine Aufgabe für ein Skript.

\begin{aufgabeBox}
Schreibe \texttt{scripts/start\_roboter.sh}, das den gesamten
Stack mit einem einzigen Befehl startet.
\end{aufgabeBox}

\begin{problemBox}
Shell-Skripte haben eine eigene Syntax. Leerzeichen um
\texttt{=} sind verboten. Zeilenenden müssen Unix-Format
haben (\texttt{\textbackslash n}, nicht \texttt{\textbackslash r\textbackslash n}).
\end{problemBox}

\begin{haubeBox}
\textbf{Warum kann man Leerzeichen um = nicht verwenden?}

In Bash ist \texttt{NAME=Wert} eine Variablenzuweisung.
\texttt{NAME = Wert} hingegen ist ein Befehl: führe das
Programm \texttt{NAME} mit den Argumenten \texttt{=}
und \texttt{Wert} aus -- und schlägt fehl, weil kein
Programm so heißt.

Das verwirrt Anfänger regelmäßig.
Die Regel: Bei Bash-Variablen niemals Leerzeichen um \texttt{=}.
\end{haubeBox}

\section{Das Startskript}

\begin{lstlisting}[style=bash]
#!/bin/bash
# scripts/start_roboter.sh
# Startet den robotik-uno-q Stack

set -e   # Skript stoppt bei erstem Fehler

PROJEKT="$HOME/robotik-uno-q"
LOG_DIR="$PROJEKT/stufe2_ros2/logs"
TIMESTAMP=$(date +"%Y%m%d_%H%M%S")

echo "=== robotik-uno-q Start ==="
echo "Zeitstempel: $TIMESTAMP"

# Verzeichnis pruefen
if [ ! -d "$PROJEKT" ]; then
    echo "FEHLER: Projektordner nicht gefunden: $PROJEKT"
    exit 1
fi

# Logverzeichnis anlegen
mkdir -p "$LOG_DIR"

# Pyserial pruefen
if ! python3 -c "import serial" 2>/dev/null; then
    echo "Installiere pyserial..."
    pip3 install pyserial
fi

# Sensor-Bridge starten
echo "Starte sensor_bridge.py..."
cd "$PROJEKT/stufe2_ros2"
python3 src/sensor_bridge.py >> "$LOG_DIR/sensor_$TIMESTAMP.log" 2>&1 &
echo "Gestartet mit PID $!"

echo "=== Stack bereit ==="
\end{lstlisting}

\begin{loesungBox}
\begin{lstlisting}[style=bash]
cd ~/robotik-uno-q
nano scripts/start_roboter.sh   # Inhalt einfuegen
chmod +x scripts/start_roboter.sh
./scripts/start_roboter.sh
\end{lstlisting}
\end{loesungBox}

\begin{projektBox}
\textbf{Bisher:} YAML-Konfiguration erstellt.\\
\textbf{Heute:} \texttt{scripts/start\_roboter.sh} -- ein Befehl
startet den gesamten Stack und schreibt Logs.\\
\textbf{Nächstes Kapitel:} Automatischer Start beim Hochfahren.
\end{projektBox}

\begin{merksatzBox}
\texttt{set -e} oben im Skript ist Pflicht: Das Skript bricht
bei jedem Fehler sofort ab, statt stumm weiterzumachen.
\texttt{\$!} ist die PID des zuletzt gestarteten Hintergrundprozesses.
\end{merksatzBox}

\begin{robotikBox}
ROS2 Launch-Dateien sind Python-Skripte (\texttt{.launch.py}).
Sie tun dasselbe wie \texttt{start\_roboter.sh} -- aber mit
ROS2-Infrastruktur: automatische Node-Discovery, Parameter-Übergabe
und geordnetes Herunterfahren.
Das Shell-Skript ist die Vorstufe, die Launch-Datei die Weiterentwicklung.
\end{robotikBox}

% ── Kapitel 10 ────────────────────────────────────────────────
\chapter{Wir wollen ein Programm regelmäßig starten}

\kapitelziel{
Du richtest den automatischen Start des robotik-uno-q-Stacks
beim Hochfahren des Raspberry Pi ein.
}
\kapitelstruktur

\section{Situation}

Beim Schulprojekt-Vortrag soll der Roboter einfach eingesteckt
werden und automatisch loslaufen. Kein Terminal aufmachen,
kein SSH, kein manueller Start. Das ist der Unterschied
zwischen einem Labormuster und einem funktionierenden System.

\begin{aufgabeBox}
Richte einen systemd-Dienst ein, der \texttt{start\_roboter.sh}
automatisch nach dem Hochfahren ausführt.
\end{aufgabeBox}

\begin{haubeBox}
\textbf{Warum systemd statt Cron?}

Cron ist für zeitgesteuerte, regelmäßige Aufgaben.
systemd verwaltet Dienste: es startet sie in der richtigen
Reihenfolge (erst Netzwerk, dann ROS2), überwacht sie
(und startet sie bei Absturz neu) und protokolliert alles.

\texttt{@reboot} in cron würde funktionieren -- aber wenn
\texttt{start\_roboter.sh} abstürzt, ist das still.
Systemd erkennt den Absturz und kann neu starten.
Für einen autonomen Roboter ist das unverzichtbar.
\end{haubeBox}

\section{Systemd-Dienst}

\begin{lstlisting}[style=text]
# /etc/systemd/system/robotik-uno-q.service
[Unit]
Description=robotik-uno-q Sensor Stack
After=network.target
Wants=network.target

[Service]
Type=simple
User=ros
WorkingDirectory=/home/ros/robotik-uno-q
ExecStart=/home/ros/robotik-uno-q/scripts/start_roboter.sh
Restart=on-failure
RestartSec=10
StandardOutput=journal
StandardError=journal

[Install]
WantedBy=multi-user.target
\end{lstlisting}

\begin{lstlisting}[style=bash]
# Dienst einrichten
sudo cp robotik-uno-q.service /etc/systemd/system/
sudo systemctl daemon-reload
sudo systemctl enable robotik-uno-q.service
sudo systemctl start robotik-uno-q.service

# Status pruefen
sudo systemctl status robotik-uno-q.service
journalctl -u robotik-uno-q.service -f   # Live-Logs
\end{lstlisting}

\begin{projektBox}
\textbf{Bisher:} Manueller Start per Skript.\\
\textbf{Heute:} Automatischer Start per systemd -- der Roboter
startet nach dem Einschalten selbstständig.\\
\textbf{Nächstes Kapitel:} Entwicklungsumgebung mit VS Code.
\end{projektBox}

\begin{merksatzBox}
\texttt{Restart=on-failure} ist der Watchdog des armen Mannes.
Wenn der Sensor-Stack abstürzt, startet systemd ihn nach
10 Sekunden neu -- ohne menschliches Eingreifen.
Das ist der Unterschied zwischen einem Prototyp und
einem produktiven System.
\end{merksatzBox}

\begin{robotikBox}
In professionellen Robotik-Systemen läuft der komplette ROS2-Stack
als systemd-Dienst. Das \texttt{docker-compose up} aus Teil~6
wird dort ebenfalls als systemd-Einheit gestartet -- dann
startet beim Einschalten automatisch Docker, dann der
ROS2-Container, dann alle Nodes.
\end{robotikBox}
